\subsection{Przebieg tras i przedmiar}
\fig{routes.pdf}{0.96}{Trasy R1--R4 prowadzone różnymi stronami obszaru. Długości obliczono z odcinków łamanych w lokalnym układzie metrycznym, bez pionów i zapasów.}
\begin{tabularx}{\linewidth}{@{}P{18mm}Y r r@{}}
\toprule\textbf{Trasa} & \textbf{Przebieg / charakter} & \textbf{Rzut [m]} & \textbf{Kabel +10\% [m]}\\\midrule
R1 & N0--(100,120)--N1; zachodnia & 320 & 352\\
R2 & N1--(100,340)--(470,340)--N2; północna & 550 & 605\\
R3 & N2--(560,240)--(560,40)--N3; wschodnia & 330 & 363\\
R4 & N3--(280,40)--N0; południowa & 320 & 352\\\midrule
\textbf{Razem} & Zakup światłowodu zaokrąglono do 1700 m. & \textbf{1520} & \textbf{1672}\\\bottomrule
\end{tabularx}

Przyjęto kanalizację dwururową: 2 $\times$ 1520 m = 3040 m, zakup \textbf{3350 m} rur z zapasem 10\%. Osiem studni stanowi założenie do weryfikacji promieni gięcia, odcinków zaciągowych i warunków gruntu. Dla czterech szaf potrzeba czterech przełącznic obsługujących po 24 włókna zakończone: łącznie 96 pigtaili LC i spawów.

Instalacje energetyczne prowadzi się w odrębnej, uzgodnionej infrastrukturze, z wymaganymi odstępami i ochroną mechaniczną. Przekroczenia dróg, głębokość, kolizje, uszczelnienia i przepusty ppoż. wymagają dokumentacji wykonawczej. Odejścia i powroty pętli SSP należy separować odpowiednio do wymaganego zachowania podczas uszkodzenia.
\begin{note}{Czego nie dowodzi rysunek}
Rozdzielenie stron pierścienia nie zapewnia pełnej niezależności: wspólne są węzły, lokalizacja N0 i część infrastruktury budynku. Trasy nie powstały na mapie do celów projektowych i nie są podstawą prowadzenia robót ziemnych.
\end{note}
